\begin{figure*}[t]
\centering
\begin{tikzpicture}[
  x=1cm,
  y=1cm,
  >={Latex[length=2mm]},
  dependency/.style={->,semithick,draw=black!65},
  module/.style={
    draw=black!65,
    rounded corners=2pt,
    align=center,
    text width=3.0cm,
    minimum height=1.35cm,
    inner sep=3pt,
    font=\footnotesize
  },
  interface/.style={module,fill=orange!10},
  application/.style={module,fill=green!10},
  reusable/.style={module,text width=3.35cm,fill=blue!8},
  result/.style={module,fill=violet!12,draw=black!80,very thick},
  lane/.style={font=\sffamily\footnotesize\bfseries,anchor=west}
]

% Application-specific path.
\node[lane] at (0,5.45) {APPLICATION-SPECIFIC SECURITY PROOF};

\node[interface] (interfaces) at (1.55,4.35) {%
  \textbf{Scheme and security interfaces}\\[0.2ex]
  \(S\), charts, correctness, and games};

\node[application] (games) at (5.15,4.35) {%
  \textbf{Noise-flooded construction and games}};

\node[application] (local) at (8.75,4.35) {%
  \textbf{One-call Gaussian bound and table invariant}};

\node[application] (adaptive) at (12.35,4.35) {%
  \textbf{Exact operation bridges and adaptive lossy hop}};

\node[result] (final) at (15.95,4.35) {%
  \textbf{Game reduction and checked theorem}\\[0.2ex]
  \(\beta_{\mathsf{CPA}}(\mathcal B)+
    \sqrt{qn}/(2\gamma)\)};

\draw[dependency] (interfaces) -- (games);
\draw[dependency] (games) -- (local);
\draw[dependency] (local) -- (adaptive);
\draw[dependency] (adaptive) -- (final);

% Reusable mathematical and program-logic path.
\node[lane] at (0,2.55) {REUSABLE FOUNDATIONS};

\node[reusable] (kl) at (1.85,1.35) {%
  \textbf{KL foundations}\\[0.2ex]
  Pinsker and the conditional-coordinate chain bound};

\node[reusable] (gaussian) at (6.15,1.35) {%
  \textbf{Discrete-Gaussian analysis}\\[0.2ex]
  Normalization, moments, and exact KL};

\node[reusable] (logic) at (10.45,1.35) {%
  \textbf{Semantic program logics}\\[0.2ex]
  Additive-error and Pythagorean rules};

\node[reusable] (compiler) at (14.75,1.35) {%
  \textbf{Selected-call compiler}\\[0.2ex]
  Adaptive replacement\\
  theorem};

\draw[dependency] (kl) -- (gaussian);
\draw[dependency] (gaussian) -- (logic);
\draw[dependency] (logic) -- (compiler);

% The two points at which the reusable layer discharges the application.
\draw[dependency] (gaussian.north) -- (local.south);
\draw[dependency] (logic.north) -- (local.south);
\draw[dependency] (compiler.north) -- (adaptive.south);

\end{tikzpicture}
\caption{Theory map of the Rocq development.  Each box is a conceptually
significant module or a small, tightly coupled module group; an arrow points
from a dependency to code that uses it.  The upper lane specializes the
reusable lower lane to noise flooding and culminates in the exported security
theorem.  Orange marks input interfaces, green the application-specific
security proof, blue reusable foundations, and purple the exported result.
Utility and adapter modules, repeated transitive imports, and proof-engineering
dependencies are omitted.}
\label{fig:theory-map}
\end{figure*}
